\section{Related Work}\label{sec:related}

\myparagraph{Learning Generative Models of the World.} Models trained to generate environment rewards and dynamics that can serve as ``world models'' for model-based reinforcement learning and planning have been recently scaled to large-scale architectures developed for vision and language ~\citep{hafner2020mastering,janner2021offline,zhang2021autoregressive,micheli2022transformers}. These works separate learning the world model from planning and policy learning, and arguably present a mismatch between the generative modeling objective of the world and learning optimal policies. Additionally, learning a world model requires the training data to be in a strict state-action-reward format, which is incompatible with the largely available datasets on the internet, such as YouTube videos. While methods such as VPT~\citep{baker2022video} can utilize internet-scale data through learning an inverse dynamics model to label unlabled videos, an inverse dynamics model itself does not support model-based planning or reinforcement learning to further improve learned policies beyond imitation learning. Our text-conditioned video policies can be seen as jointly learning the world model and conducting hierarchical planning simultaneously, and is able to leverage widely available datasets that are not specifically designed for sequential decision making. 

\myparagraph{Diffusion Models for Decision Making.} 
Diffusion models have recently been applied to different decision making problems~\citep{janner2022planning,ajay2022conditional,wang2022diffusion, zhang2022lad, urain2022se, zhong2022guided}. 
Most similar to this work, \citet{janner2022planning} trained an unconditional diffusion model to generate trajectories consisting of joint-based states and actions, and used a separately trained reward model to select generated plans. On the other hand, \citet{ajay2022conditional} %
trained a conditional diffusion model to guide behavior synthesis from desired rewards, constraints or agent skills. Unlike both works, which learn task-specific policies from scratch, our approach of text-condition video generation as a universal policy can leverage internet-scale knowledge to learn generalist agents that can be deployed to a variety of novel tasks and environments. Additionally, \citet{kapelyukh2022dall} applied web-scale text-conditioned image diffusion to generate a goal image to condition a policy on, whereas our work uses video diffusion to learning universal policies directly.

\myparagraph{Learning Generalist Agents.} Inspired by the success of large-scale pretraining in vision and language domains, large-scale sequence and image models have recently been applied to learning generalist decision making agents~\citep{reed2022generalist,lee2022multi,kumar2022offline}. However, these generalist agents can only operate under environments with the same state and action spaces (e.g., Atari games)~\cite{lee2022multi,kumar2022offline}, or require studious tokenization~\citep{reed2022generalist} that might seem unnatural in %
scenarios
where different environments have distinct state and actions spaces. Another downside of using customized tokens for control is the inability to directly utilize knowledge from pretrained vision and language models. Our approach, on the other hand, uses text and images as universal interfaces for policy learning so that the knowledge from pretrained vision and language models can be preserved. The choice of diffusion as opposed to autoregressive sequence modeling also enables long-term and hierarchical planning. 

\paragraph{Learning Text-Conditioned Policies.} There has been a growing amount of work using text commands as a way to learn multi-task and generalist control policies~\citep{huang2022inner,ahn2022can,brohan2022rt,lynch2020language,mahmoudieh2022zero,liu2022instruction}. Different from our framing of video-as-policies, existing work directly trains a language-conditioned control policy in the action space of some specific robot, leaving cross-morphology multi-environment learning of generalist agents as an unsolved problem. We believe this paper is the first to propose images as a universal state and action space to enable broad knowledge transfer across environments, tasks, and even between humans and robots.
